\documentclass[aspectratio=169]{beamer}
\usepackage[%
numbering=fraction,%
block=fill,%
subsectionpage=progressbar]{pescPres}
\usepackage{appendixnumberbeamer}
\usepackage[brazilian]{babel}
\usepackage{booktabs}

\title[Metodologia Experimental]{Metodologia Experimental}
\subtitle{Resumo para apresentação de 20 minutos \\ v-2.0 (revisão crítica)}
\author{João Pedro Vital Brasil Wieland}
\advisor{Claudia Maria Lima Werner}

\begin{document}
\maketitle

%================================================
\section{}

\begin{frame}{O problema e a mensagem central}
\small
\begin{itemize}
  \item \textbf{Problema}: correção automática de acessibilidade em aplicações React é difícil porque:
  \begin{itemize}
    \item scanners têm falsos positivos/negativos,
    \item violações variam muito em complexidade,
    \item correções podem \textbf{quebrar} funcionalidade.
  \end{itemize}
  \item \textbf{Tese metodológica}: para avaliar reparo com LLMs de forma científica, precisamos de:
  \begin{itemize}
    \item \textbf{realismo} (projetos reais),
    \item \textbf{controle} (condições comparáveis),
    \item \textbf{reprodutibilidade} (snapshot + ordem determinística),
    \item \textbf{segurança} (validação anti-regressão).
  \end{itemize}
  \item \textbf{Contribuição}: um pipeline reprodutível e controlado para
        \textbf{detectar, corrigir e validar} violações WCAG com LLMs.
\end{itemize}
\end{frame}

\begin{frame}{Por que esta metodologia é ``a mais científica'' para o problema}
\small
\begin{itemize}
  \item \textbf{Evita vieses clássicos} de avaliações com LLMs:
  \begin{itemize}
    \item exemplos sintéticos e fáceis,
    \item ordem aleatória de processamento,
    \item sucesso sem checar regressão,
    \item comparações injustas (problemas diferentes por condição).
  \end{itemize}
  \item \textbf{Controla variáveis críticas}:
  \begin{itemize}
    \item mesmos problemas para todas as condições (comparação pareada),
    \item parâmetros fixos do modelo (redução de aleatoriedade),
    \item validação padronizada (mesmo oráculo).
  \end{itemize}
  \item \textbf{É honesta com suas limitações}: ameaças à validade são
        \textbf{documentadas e mitigadas} no desenho, não ignoradas.
  \item \textbf{Torna o resultado auditável}:
  cada correção é um patch rastreável;
  cada decisão é operacionalizada (equações).
\end{itemize}
\end{frame}

%================================================
\section{Ferramenta Desenvolvida}

\begin{frame}{A ferramenta: o que ela faz (visão prática)}
\small
\begin{itemize}
  \item A ferramenta implementa uma \textbf{linha de produção}:
  \begin{enumerate}
    \item baixa e congela projetos React (\emph{snapshot} por commit);
    \item executa múltiplos scanners e normaliza violações;
    \item prioriza violações de forma determinística;
    \item escolhe estratégia/agente com base em complexidade (\emph{routing});
    \item gera patch com LLM (prompts controlados);
    \item valida o patch com análise estática em 4 camadas;
    \item reexecuta scanners e registra métricas.
  \end{enumerate}
  \item Resultado final: um conjunto de \textbf{correções validadas} +
        um log completo para avaliação científica.
\end{itemize}
\end{frame}

\begin{frame}{Por que multiagente (e não um único prompt gigante)?}
\small
\begin{itemize}
  \item \textbf{Observação prática}: violações WCAG diferem muito em escopo e
        complexidade.
  \item Um único agente tende a:
  \begin{itemize}
    \item desperdiçar contexto em casos simples,
    \item falhar em casos que exigem vários arquivos,
    \item gerar patches grandes e arriscados.
  \end{itemize}
  \item A arquitetura multiagente permite:
  \begin{itemize}
    \item agentes especializados por tipo de tarefa (local vs.\ global),
    \item decisões explícitas sobre contexto e restrições,
    \item controle fino para ablações e comparação justa.
  \end{itemize}
  \item \textbf{Roteamento explícito}: a decisão entre agentes é registrada
        com \emph{score} e razão — cada escolha é auditável.
\end{itemize}
\end{frame}

%================================================
\section{Escolhas Experimentais (e por quê)}

\begin{frame}{Escolha 1: projetos reais + snapshot por commit}
\small
\begin{itemize}
  \item \textbf{Por que projetos reais?}
  \begin{itemize}
    \item evita ``toy problems'',
    \item incorpora build, dependências, estilos e dívida técnica reais.
  \end{itemize}
  \item \textbf{Por que snapshot?}
  \begin{itemize}
    \item garante que o código não muda durante o experimento,
    \item hash SHA-256 do conteúdo protege contra mudanças silenciosas,
    \item permite reprodução exata no futuro.
  \end{itemize}
  \item \textbf{Efeito científico}: aumenta validade externa sem sacrificar
        reprodutibilidade.

  \medskip
  \alert{\textbf{Ameaça reconhecida}}: componentes são escaneados em
  \emph{harness} HTML isolado, sem contexto de routing ou estado global.
  Violações dependentes de contexto podem ser sub/super-estimadas.
  $\Rightarrow$ Escopo declarado: desempenho ao nível do scanner
  (\emph{scanner-level}), não ao nível da aplicação integrada.
\end{itemize}
\end{frame}

\begin{frame}{Escolha 2: múltiplos scanners + normalização}
\small
\begin{itemize}
  \item \textbf{Por que múltiplos scanners?}
  \begin{itemize}
    \item cada scanner tem heurísticas e coberturas diferentes,
    \item reduz dependência de um único oráculo imperfeito,
    \item aumenta \emph{recall} do conjunto de violações detectadas.
  \end{itemize}
  \item \textbf{Confiança baseada em consenso}:
  \begin{itemize}
    \item \textbf{HIGH}: $\geq$ \texttt{min\_tool\_consensus} ferramentas concordam
          no mesmo seletor + critério WCAG,
    \item \textbf{MEDIUM}: 1 ferramenta + impacto \emph{critical/serious},
    \item \textbf{LOW}: demais casos.
  \end{itemize}
  \item \textbf{Threshold configurável} (\texttt{.env}): permite ajustar o
        critério de consenso sem mudar código.
  \item \textbf{Por que normalizar?} Saídas são heterogêneas (formatos e
        severidades); normalização cria a unidade experimental estável:
        \textbf{``violação tratável''}.
  \item \textbf{Efeito científico}: aumenta controle e comparabilidade entre execuções.
\end{itemize}
\end{frame}

\begin{frame}{Escolha 3: ordem determinística via função de sorting}
\small
\begin{itemize}
  \item \textbf{Problema}: ordem aleatória altera consumo de recursos e resultados.
  \item \textbf{Solução}: ordenar violações por uma chave determinística e justificável.
\end{itemize}
\vspace{0.4em}
\[
  \text{sort\_key}(i)=\bigl(
  -\text{conf\_rank}(i.\texttt{confidence}),\;
  -\text{imp\_rank}(i.\texttt{impact}),\;
  i.\texttt{wcag\_criteria},\;
  i.\texttt{selector}
  \bigr)
\]
\vspace{0.2em}
\begin{itemize}
  \item \textbf{Justificativa}: corrige primeiro o que é mais confiável
        e mais severo.
  \item \textbf{Sentinel}: critérios sem mapeamento WCAG recebem
        \texttt{"9.9.9"} $\Rightarrow$ ordenados ao final
        (decisão explícita, não implícita).
  \item \textbf{Efeito científico}: elimina um confound importante
        (ordem de processamento).
\end{itemize}
\end{frame}

\begin{frame}{Escolha 4: comparação ``justa'' (mesmos problemas por condição)}
\small
\begin{itemize}
  \item \textbf{Princípio}: uma condição não pode receber um conjunto mais fácil de violações.
  \item \textbf{Desenho}: cada violação é tentada em múltiplas condições
        (modelo/prompt/ablação).
  \item Consequências:
  \begin{itemize}
    \item comparações são \textbf{pareadas} (\emph{within-subject}),
    \item variação por dificuldade do problema é controlada,
    \item diferenças observadas são mais atribuíveis à condição.
  \end{itemize}
  \item \textbf{Testes estatísticos compatíveis}:
  \begin{itemize}
    \item H1: McNemar (pares binários),
    \item H2--H3: Kruskal-Wallis + Mann-Whitney pós-hoc + Bonferroni.
  \end{itemize}
  \item \textbf{Efeito científico}: aumenta poder estatístico e reduz viés de seleção.
\end{itemize}
\end{frame}

\begin{frame}{Escolha 5: dois estudos de ablação para explicar ``o que importa''}
\small
\begin{columns}[T]
\column{0.48\textwidth}
\textbf{Ablação do template de prompt} (H1 --- confirmatory)
\begin{itemize}
  \item Template tem 6 componentes:\\
        role, constraints, file, issues, few-shot examples, output format
  \item Remove um componente por vez e mede SR
  \item McNemar + Cliff's $\delta$
  \item Responde: quais partes são essenciais?
\end{itemize}

\column{0.48\textwidth}
\textbf{Ablação de scanners} (exploratório)
\begin{itemize}
  \item Variantes: \emph{pa11y\_only}, \emph{axe\_only},
        \emph{playwright\_only}, \emph{pa11y+axe}, \emph{all\_tools}
  \item Modelo fixo (isola variável dos scanners)
  \item Mede SR, IFR, MTTR por combinação
  \item Responde: quanto cada ferramenta contribui?
\end{itemize}
\end{columns}

\vspace{0.5em}
\small
\alert{Distinção fundamental}: ablação de prompt é \textbf{confirmatória}
(H1 pré-especificada); ablação de scanners é \textbf{exploratória}
(gera hipóteses para trabalhos futuros).
\end{frame}

%================================================
\section{Segurança e Validação}

\begin{frame}{Escolha 6: validação em 4 camadas (análise estática)}
\small
\begin{itemize}
  \item \textbf{Princípio}: ``corrigir sem quebrar'' é requisito do mundo real.
  \item Um patch só é aceito se passar por quatro camadas de análise estática:
  \begin{enumerate}
    \item \textbf{Sintática}: heurísticas de formato (resposta não vazia,
          bloco JSX presente, sem marcadores de recusa do LLM),
    \item \textbf{Preservação estrutural}: interfaces TypeScript, assinaturas de
          export e handlers de evento preservados por análise de texto,
    \item \textbf{Verificação de domínio}: presença dos atributos de
          acessibilidade esperados pelos tipos de violação corrigidos,
    \item \textbf{Qualidade de código}: padrões proibidos
          (\texttt{tabIndex~$<$~-1}, \texttt{dangerouslySetInnerHTML}).
  \end{enumerate}
  \item \textbf{Escopo declarado}: as camadas são \emph{análise estática},
        não execução de runtime.
        Uma aprovação \textbf{não garante} ausência de regressão de
        comportamento --- isso é uma ameaça documentada (Seção~3.8.1).
\end{itemize}
\end{frame}

\begin{frame}{Métricas: eficácia e segurança, sem ``truques''}
\small
\begin{itemize}
  \item \textbf{Sucesso} é definido de forma verificável e operacional.
  \item Métricas centrais:
\end{itemize}
\vspace{0.2em}
\[
  \text{SR} =
  \frac{\#\text{arquivos que passaram todas as 4 camadas}}{\#\text{tentativas}}
  \qquad
  \text{IFR} =
  \frac{\sum|I_{\text{corrigidas}}(f)|}{\sum|I(f)|}
\]
\[
  \text{MTTR} =
  \frac{\sum t_{\text{corrigido}}}{\#\text{sucessos}}
  \qquad
  \text{TE} =
  \frac{\text{IFR} \times |I|}{\text{tokens}/1000}
\]
\vspace{0.2em}
\begin{itemize}
  \item \textbf{Inferência} vai além do valor-$p$:
  \begin{itemize}
    \item critério duplo: $p < \alpha$ \textbf{e} $|\delta| \geq 0.147$
          (efeito pequeno de Cliff),
    \item bootstrap CI com $n=2000$, semente 42 (reprodutível),
    \item Bonferroni para comparações múltiplas (H2--H3).
  \end{itemize}
  \item \textbf{MTTR}: medida sobre sucessos apenas (MTTR indefinida em
        zero sucessos --- não é um ``truque'' estatístico, é o comportamento
        correto da métrica).
\end{itemize}
\end{frame}

%================================================
\section{Ameaças à Validade}

\begin{frame}{Ameaças à validade (reconhecidas e mitigadas)}
\small
\begin{columns}[T]
\column{0.48\textwidth}
\textbf{Validade de construto}
\begin{itemize}
  \item Scanners em \emph{harness} isolado $\neq$ app integrada\\
        \textit{Mitigação}: escopo declarado como \emph{scanner-level}
  \item Validação estática $\neq$ ausência de regressão de runtime\\
        \textit{Mitigação}: $\rho$ (H5) mede taxa de rejeições --- limitação documentada
\end{itemize}

\textbf{Validade interna}
\begin{itemize}
  \item Threshold do router calibrado em 50 arquivos\\
        \textit{Mitigação}: comentário no código + seção de limitações
  \item LLM não determinístico por natureza\\
        \textit{Mitigação}: temperatura=0.1, 3+ repetições por condição
\end{itemize}

\column{0.48\textwidth}
\textbf{Validade externa}
\begin{itemize}
  \item Apenas modelos locais (Qwen, DeepSeek, CodeLlama\ldots)\\
        \textit{Mitigação}: resultados são interpretados para esse escopo;
        generalização para modelos proprietários é trabalho futuro
  \item Projetos selecionados por disponibilidade pública
\end{itemize}

\textbf{Validade de conclusão estatística}
\begin{itemize}
  \item Múltiplas comparações\\
        \textit{Mitigação}: Bonferroni
  \item Pequenas amostras por categoria\\
        \textit{Mitigação}: bootstrap CI + Cliff's $\delta$
\end{itemize}
\end{columns}
\end{frame}

%================================================
\section{Fechamento}

\begin{frame}{Por que confiar nos resultados}
\small
\begin{itemize}
  \item O experimento é confiável porque:
  \begin{itemize}
    \item usa projetos reais e congelados por hash (reprodutível),
    \item controla ordem e condições (determinístico),
    \item compara de forma pareada (justo),
    \item documenta o que a validação faz e o que não faz (honesto),
    \item mede com métricas simples e auditáveis.
  \end{itemize}
  \item \textbf{Interpretação correta}: resultados indicam capacidade de
        reparo \emph{scanner-level} sob validação estática anti-regressão.
  \item \textbf{Distinção confirmatory/exploratory} previne
        \emph{p-hacking}: H1--H4 pré-especificadas; H5 gera
        hipóteses para trabalhos futuros.
  \item \textbf{O que isso entrega}: um protocolo científico replicável para
        avaliação de reparo automático de acessibilidade com LLMs locais.
\end{itemize}
\end{frame}

\appendix
\begin{frame}[allowframebreaks]
  \frametitle{Referências}
  \nocite{*}
  \bibliographystyle{apalike}
  \bibliography{biblio.bib}
\end{frame}

% ── Slides extras (backup para perguntas da banca) ──────────────────────────

\begin{frame}[noframenumbering]{Backup: como o router decide o agente}
\small
\begin{columns}[T]
\column{0.55\textwidth}
\textbf{Matriz de pontuação} (código: \texttt{router/engine.py})
\begin{itemize}
  \item[\textbf{+4}] tipos CONTRAST ou SEMANTIC presentes
  \item[\textbf{+4}] $n_{\text{issues}} \geq \tau$ (\texttt{swe\_max\_issues}=4)
  \item[\textbf{+5}] $n_{\text{issues}} \geq 2\tau$
  \item[\textbf{+3}] algum issue com complexidade COMPLEX
  \item[\textbf{+3}] $\geq$ 3 tipos de issue distintos
  \item[\textbf{-3}] todos ARIA/LABEL/ALT-TEXT e $n<\tau$
\end{itemize}
\vspace{0.3em}
$\text{score} \geq 3 \Rightarrow$ OpenHands \quad
$\text{score} < 3 \Rightarrow$ SWE-agent

\column{0.42\textwidth}
\textbf{Limitação declarada} (\texttt{ROUTING\_THRESHOLD = 3}):
\begin{quote}\small
\textit{``empirically calibrated on 50 representative files. Internal validity threat: if the benchmark distribution differs from the calibration set, this threshold may be suboptimal.''}
\end{quote}
$\Rightarrow$ Threshold é uma \textbf{decisão de design
explicitamente documentada}, não um fato empírico validado.
\end{columns}
\end{frame}

\begin{frame}[noframenumbering]{Backup: o que a validação realmente verifica}
\small
\begin{tabular}{@{}llp{6cm}@{}}
\toprule
\textbf{Camada} & \textbf{Técnica} & \textbf{O que verifica / O que não verifica} \\
\midrule
1 & Heurística regex & Não vazio, bloco JSX presente, sem recusa LLM \\
  &                  & \alert{Não} é parse AST completo \\
\addlinespace
2 & Análise estática & Interfaces TS, exports e event handlers preservados \\
  &                  & \alert{Não} compila, \alert{não} roda testes \\
\addlinespace
3 & Heurística regex & Atributos \texttt{alt=}, \texttt{aria-label=} esperados \\
  &                  & \alert{Não} re-executa os scanners reais \\
\addlinespace
4 & Pattern matching & \texttt{tabIndex < -1}, \texttt{dangerouslySetInnerHTML} \\
\bottomrule
\end{tabular}

\vspace{0.5em}
\small
\textbf{Implicação}: aprovação em todas as camadas é condição necessária
mas não suficiente para ausência de regressão de runtime.
A \textbf{taxa de regressão} $\rho$ (H5, exploratória) mede exatamente
o que a camada~2 detecta estaticamente --- e é apresentada como
\emph{exploratory}, não como prova definitiva.
\end{frame}

\begin{frame}[noframenumbering]{Backup: confirmatory vs.\ exploratory}
\small
\begin{columns}[T]
\column{0.48\textwidth}
\textbf{Confirmatory (H1--H4)} --- pré-especificadas
\begin{itemize}
  \item H1: ablação de prompt melhora SR?
        (McNemar + $|\delta| \geq 0.147$)
  \item H2: qual estratégia de prompt é melhor?
        (Kruskal-Wallis + TE $\geq$ 80\% zero-shot)
  \item H3: arquitetura do LLM importa?
        (Kruskal-Wallis pairwise)
  \item H4: complexidade da violação prediz sucesso?
        (direcional: IFR\textsubscript{simples} $>$ IFR\textsubscript{complexo})
\end{itemize}

\column{0.48\textwidth}
\textbf{Exploratory (H5 + outros)}
\begin{itemize}
  \item H5: taxa de regressão $\rho$ por (modelo, categoria)\\
        com bootstrap CI --- \emph{descritivo}
  \item Ablação de scanners: qual combinação de ferramentas é mais eficaz?
  \item Perfil de rejeição por camada
\end{itemize}

\vspace{0.5em}
\alert{Por que isso importa?}\\
A distinção previne \emph{p-hacking} post-hoc:
resultados exploratórios geram hipóteses para \emph{futuras} avaliações,
não confirmam as hipóteses do estudo corrente.
\end{columns}
\end{frame}

\end{document}
